\clearpage\thispagestyle{empty}
\pdfbookmark[0]{Sources for Gauss to Bode (7 September 1818)}{bodesep1818-sources}
\begin{center}\textbf{Sources for Gauss to Bode (7 September 1818)}\end{center}
\noindent This AI translation follows Carl Friedrich Gauss’s letter to Johann Elert Bode in \href{https://gdz.sub.uni-goettingen.de/id/PPN236007467}{\textit{Werke}, vol.~VI} (Göttingen, 1874), pp.~414--416. It first appeared under the title \textit{Ueber den Repsoldschen Meridiankreis, Beobachtung des Uranus, Saturns und der Pallas etc.} as item~24 in Bode’s \href{https://api.digitale-sammlungen.de/iiif/image/v2/bsb10538350_00003/full/1600,/0/default.jpg}{\textit{Astronomisches Jahrbuch für das Jahr 1821}} (Berlin, 1818), pp.~\href{https://api.digitale-sammlungen.de/iiif/image/v2/bsb10538350_00218/full/1600,/0/default.jpg}{212}--\href{https://api.digitale-sammlungen.de/iiif/image/v2/bsb10538350_00222/full/1600,/0/default.jpg}{216}. The yearbook links lead to the Bavarian State Library copy.
\par\smallskip\noindent Two notes appear on p.~\href{https://api.digitale-sammlungen.de/iiif/image/v2/bsb10538350_00219/full/1600,/0/default.jpg}{213} of the first printing. The first is attached to the mention of the Pallas ephemeris: “Follows later.” The ephemeris calculated by Dirksen appears separately as item~30, pp.~\href{https://api.digitale-sammlungen.de/iiif/image/v2/bsb10538350_00230/full/1600,/0/default.jpg}{224}--\href{https://api.digitale-sammlungen.de/iiif/image/v2/bsb10538350_00231/full/1600,/0/default.jpg}{225}, and covers 7 August 1819 to 6 June 1820 for midnight in Göttingen. The second note is attached to Gauss’s remark that Bode had doubtless seen the circle: “The ingenious instrument maker showed it to me when I was in Hamburg in 1810. B.”
\par\smallskip\noindent In the first printing, p.~\href{https://api.digitale-sammlungen.de/iiif/image/v2/bsb10538350_00220/full/1600,/0/default.jpg}{214} labels the first group of Polaris observations \textit{Jul.}, where the collected text has \textit{Juni}. On p.~\href{https://api.digitale-sammlungen.de/iiif/image/v2/bsb10538350_00221/full/1600,/0/default.jpg}{215} the mean correction is marked with a double prime, where the collected text has a raised \textit{s}. The Uranus declination on 3 June ends in $41'',5$ in the first printing and $41''6$ in the collected letter. Both the irregular verb after the singular “column” and the spelling “obsorvations” reflect the collected text. Translation and typesetting produced by AI.
\clearpage
